% yugen-ash highlight showcase — LaTeX (.tex)
%
% Open with the theme:
%   nvim -u tests/minimal_init.lua tests/fixtures/langs/latex.tex
%   YUGEN_ASH_VARIANT=light nvim -u tests/minimal_init.lua tests/fixtures/langs/latex.tex
%
% Headless coverage check:
%   scripts/test-highlights.sh latex
%
% Debug: :Inspect, :InspectTree, :YugenAshInspect, :hi <group>

% legacy syntax: texComment
% Treesitter: @comment
% Plain comment for texComment.
% TODO: exercise texTodo with a real TODO marker.

% legacy syntax: texDocType, texDocTypeArgs
% Treesitter: @keyword
\documentclass[11pt,a4paper]{article}

% legacy syntax: texStatement, texCmdName, texCmdArgs, texOption
% Treesitter: @function, @function.macro, @punctuation.delimiter
\usepackage[utf8]{inputenc}
\usepackage{amsmath,amssymb}
\usepackage[colorlinks=true,linkcolor=blue]{hyperref}

% legacy syntax: texNewCmd, texDefName, texCmdArgs
% Treesitter: @function.macro, @keyword
\newcommand{\showcase}[1]{\textbf{#1}}
\newenvironment{note}[1]{\begin{quote}\itshape #1}{\end{quote}}

% legacy syntax: texInputFile
\input{sections/intro}

\begin{document}

% legacy syntax: texSection, texSectionMarker, texSectionName
% Treesitter: @markup.heading
\section{Introduction}\label{sec:intro}
\subsection{Scope}
\subsubsection{Details}

% legacy syntax: texBeginEnd, texBeginEndName
% Treesitter: @markup.environment, @markup.environment.name
\begin{note}{A note}
	This is a \showcase{demonstration} of the theme.
\end{note}

% legacy syntax: texRefZone, texRefLabel, texLabel, texCite
% Treesitter: @markup.link, @markup.link.label
See Section~\ref{sec:intro} and \cite{knuth1984texbook}.
Refer to \label{sec:details}.

% legacy syntax: texType, texTypeStyle
% Treesitter: @type
Text can be \textbf{bold}, \textit{italic}, \texttt{monospace}, and \emph{emphasised}.

% legacy syntax: texMathZone, texMath, texMathSymbol, texMathOper
% legacy syntax: texMathMatcher, texMathDelim, texMathDelimZone
% legacy syntax: texGreek, texSuperscript, texSubscript
% Treesitter: @markup.math, @punctuation.special
The equation
\[
	\alpha^2 + \beta_1^2 = \gamma_{max} \sum_{i=1}^{n} x_i
\]
is inline math: $E = mc^2$ and \(\theta \leq \pi\).

% legacy syntax: texString, texSpecialChar
% Treesitter: @string
\url{https://example.com/~showcase}
A percent sign: \%, ampersand: \&, underscore: \_, hash: \#.

% legacy syntax: texDelimiter, texOption
% Treesitter: @punctuation.delimiter
\textcolor{red}{Highlighted text}

% legacy syntax: texMathMatcher
\begin{align}
	a &= b + c \\
	d &= e - f
\end{align}

% legacy syntax: texBeginEnd, texBeginEndName
\end{document}
